\subsection{A compatible degree-seven replacement}

\begin{statement}{D.31}{Proved}{A class-three two-primary kernel with third subgroup of rank 1860}
Retain $K,N,\mathcal P_2$ and the Zassenhaus filtration
from D.30. Put
\[
 T_2=B/D_8(K),\qquad
 U=\gamma_3N\,D_8(K)/D_8(K),\qquad
 C_{**}=B/(\gamma_3N\,D_8(K)).                         \tag{D.31.1}
\]
The group $U$ is elementary abelian of dimension 1860.
Its degree-six quotient has dimension 240 and its
degree-seven subgroup dimension 1620. Conjugation
makes $U$ an actual $Q$-module. Its cusp-correction
map has rank 1810 and cokernel dimension 50.

The fixed homology and CM data of D.30 admit another
explicit central replacement. Its fifty degree-seven
conditions, together with the fifteen class-two
$\mathbb Z/4$ conditions, form a map of rank 65 on the
4769-dimensional central relation space. The resulting
4,903,248-letter cusp word vanishes in both complete
models: the rank-97 class-two Magnus model over
$\mathbb Z/4$, and the rank-four mod-two Magnus model
through degree seven. Its multiplier is 1. This gives
compatible automorphisms of $H^{(2)}$ and $T_2$, with
\[
 |T_2|=576\cdot2^{3333},\quad
 |C_{**}|=576\cdot2^{1473},\quad
 |H^{(2)}\times_{C_{**}}T_2|=576\cdot2^{11463}.
                                                               \tag{D.31.2}
\]

The same stable correction construction extends this
seed to the precise quotient
\[
 \mathcal S_2^{(7)}=
 \mathcal P_2/\bigl(
 \overline{\gamma_3N^{(2)}}\cap D_8(K^{(2)})\bigr)
 \simeq\mathcal P_2^{[2]}\times_{C_{**}}T_2.            \tag{D.31.3}
\]
Its kernel over $Q$ has class exactly three, with
third lower-central subgroup $C_2^{1860}$ and trivial
fourth subgroup. Its multiplier has two-adic component
in $1+16\mathbb Z_2$. The odd-prime components can all
be 1, with the original odd-primary conjugators retained
separately. The action combines with every odd relative
pro-prime completion and D.28's mixed quotient $J$.

There is a canonical surjection
\[
 \pi:\mathcal S_2^{(7)}\longrightarrow\mathcal S_2
                                                               \tag{D.31.4}
\]
to D.30's target, whose kernel is $C_2^{1620}$.
The recorded new automorphism induces an automorphism
$\alpha_6'$ there, with
$\pi\alpha_7=\alpha_6'\pi$.
It is an explicitly different central representative
from the recorded D.30 automorphism $\alpha_6$:
their reductions already differ on the image of $s$
in $H^{(1)}$. The nonzero central difference is
recorded in the reduction certificate. Thus this
step replaces the earlier action within its fixed
homology class; it does not assert that the earlier
central coordinates were all retained.

One global cusp gauge gives compatible nonarithmetical
finite shadows on characteristic refinements retaining
the $G_6$ and CM data. D.32 below proves that the
recorded $T_2$ action has no peripheral lift to the
full relative pro-two source, hence no genuine full
profinite lift. The attainability of the same extra
CM coset by a different action remains undecided.
\end{statement}
\begin{proof}
Modulo $D_4(K)$ every element of $N$ is a product of
the nine chosen degree-two relation lifts and elements
of $D_3(K)/D_4(K)$, whose twenty Lyndon brackets are
retained. A triple commutator modulo $D_8(K)$ uses
only degree patterns $(2,2,2)$ and permutations of
$(3,2,2)$. All non-Lie correction terms begin in degree
eight. Products of these degree-six and degree-seven
elements are additive, and their squares vanish.
Conjugation by a degree-one element of $K$ adds a
degree-seven term $[L_1,W]$. Jacobi expresses it using
the same generators, because $[L_1,R_2]\subset L_3$.
Thus the normal closure contributes no omitted
degree-seven generators.

The program retains the 240 actual triple-commutator
words chosen in D.30, including their degree-seven
terms. It adjoins the brackets
\[
       [r_i,[r_j,z_k]],\qquad
       1\leq i,j\leq9,\quad 1\leq k\leq20,            \tag{D.31.5}
\]
where the $r_i$ are the degree-two relation lifts and
the $z_k$ are the degree-three Lyndon lifts. Their
leading degree-seven terms have rank 1620. The earlier
240 basis elements have independent degree-six terms,
so the combined rank is 1860. This proves the exact
description of $U$, with its displayed filtration.
Since $[N,U]\subset[D_2(K),D_6(K)]\subset D_8(K)$,
its conjugation action factors through $Q=B/N$.

The action retains the extension between degrees six
and seven. The 240 inhomogeneous basis words are
conjugated by the actual $B$ words $s$ and $sr$ and
expanded through degree seven. Only the pure
degree-seven basis uses the leading tensor action.
The actual elliptic images may use these same module
actions: $q(\beta)$ centralizes $s$ in $Q$, so
$\beta s\beta^{-1}s^{-1}\in N$, whose conjugation on
$U$ is trivial; likewise the $A$ conjugate of $r$
differs from $r$ by an element of $N$.
This proves the precise action entering
\[
                 (1-s)b+s(1-r)a\quad(b,a\in U).        \tag{D.31.6}
\]
Its exact rank is 1810. The code obtains all fifty
annihilating functionals from this actual module.

For simultaneous correction, use $C_1$ from D.29.
Its generators have degree at least four in $K$,
so its image is abelian modulo $D_8(K)$. Its effect
in that quotient depends only on its coefficients
modulo two. In the $N$ class-two $\mathbb Z/4$ model
it is the earlier abelian module of exponent four.
The degree-seven cusp effects use the actual elliptic
conjugate path through degree three and the images of
the 97 $N$ generators through degree five. Here the
degree-two action of $N$ does matter on $C_1$; the
calculation therefore does not replace it by the
leading $Q_0$ tensor matrices.

On the 4769-dimensional relation space, impose the
fifteen conditions for lifting the correction
modulo four and the fifty functionals from (D.31.6).
The joint rank is 65. Its solved target retains zero
in the first fifteen coordinates and cancels the
new cusp obstruction in the last fifty. Functionals
on $U$ are extended linearly for evaluating individual
columns; this introduces no assumption, because the
class-two conditions put the final combined error
in $U$. An upper twice-central correction then
restores the complete $\mathbb Z/4$ equation and
lies in $D_8(K)$, so it changes no term of the other
equation. Finally the actual triple-commutator words
representing $U$ remove its error. They vanish in
every class-two quotient of $N$.

The resulting 4,903,248-letter $B$ word is evaluated
directly in both group models. All rank-97 degree-one
and degree-two coefficients modulo four vanish,
and all rank-four coefficients in degrees one
through seven modulo two vanish. Its multiplier is 1.
The retained homology map on $N^{\rm ab}/2$ has
rank 97; the action on $K/D_2(K)$ is the identity
through $O_2(Q)^{\rm ab}$. The finite two-kernel
Frattini argument proves automorphisms of both
finite quotients, in addition to their cusp equations.

The next restricted Witt dimension is
$\dim D_7(K)/D_8(K)=(4^7-4)/7=2340$, using the same
Hall-power basis in \cite[Proposition 4.4]{dEfrat2023}.
The sum through degree seven is 3338. Hence
$|T_2|=18\cdot2^{3338}=576\cdot2^{3333}$.
The kernel of $T_2\to C_{**}$ is exactly $U$,
of order $2^{1860}$. As in D.30,
\[
 (\gamma_3N[N,N]^4N^8)D_8(K)=\gamma_3N D_8(K),         \tag{D.31.7}
\]
because the last two factors lie in $D_{16}(K)$.
This proves the exact common quotient, the remaining
orders in (D.31.2), and the canonical surjection
from $B$ onto the finite fibre product. The same
two-primary conjugator words give agreement there.

Restart the stable iteration from the new seed.
Its $G_2$ action has not changed, so the stabilized
correction matrices and their fifteen-dimensional
surjectivity certificate are unchanged. Every later
$C_k$ correction, $k\geq2$, lies in $D_8(K)$, so the
entire $T_2$ action and its conjugators stay fixed.
The cusp $t$ still has order 48 in $T_2$: the
degree-one term of $t^6$ is nonzero, its fourth
power has a nonzero degree-four term, and its eighth
power starts in degree eight. The later multiplier
increments are divisible by 144 and therefore act
trivially on $T_2$. They converge to a two-adic
multiplier in $1+16\mathbb Z_2$.
The inverse-limit normal-subgroup formula gives
(D.31.3). Its third lower-central subgroup over $Q$
is $U$, and its fourth is trivial by
$[N,\gamma_3N]\subset D_8(K)$. The nonzero third
subgroup proves that the class is exactly three.

The inclusion $D_8(K)\subset D_7(K)$ gives (D.31.4).
Its kernel is the degree-seven part of the image of
$\gamma_3N$, of dimension 1620. The new automorphism
preserves this quotient: it preserves the preimages
of $N$ and $K$ over $Q$, their lower-central subgroups,
and the characteristic Zassenhaus filtration. The
comparison of recorded actions is checked separately
in \texttt{d\_cm\_degree7\_reduction.py}. Let $z$ be
the lower beta central-change vector and $M_s$ the
rank-97 homology matrix. Its associated symmetric
matrix $Z$ has nonzero difference
\[
                       Z+M_sZM_s^{\mathsf T}.          \tag{D.31.8}
\]
The twice-central and triple-commutator corrections
vanish in $H^{(1)}$, so (D.31.8) is precisely the
change of the $s$ image there. This proves
$\alpha_6'\ne\alpha_6$ by an explicit quotient map,
while the original $G_6$ homology data remain fixed.

For odd primes choose the original conjugators as
separate components, using the coprime centralizer
construction. The discrete two-primary word is not
asserted to preserve its odd-primary abelianization.
The new two-primary kernel is a pro-two group, so
the maximal common quotient with $J$ is still the
$H\times_QG_3$ proved in D.28 after adjoining the
odd components. The homology and separately chosen
three-primary data agree there. The exact fibre
product consequently carries one combined action
and compatible conjugators. The global cusp gauge
supplies the claimed finite shadow family. The
quotients (D.31.1)--(D.31.4) specify the retained
information. D.32 applies the next finite quotient
to the recorded action and proves a complete
obstruction there; the quotient constructions
proved here remain valid.
\end{proof}
